\medskip
\hypertarget{native-109-article080}{}\pdfbookmark[2]{80.}{article-080}
\begin{center}80.\end{center}
Si angulus $\Pi$ datus est, $p$ et $e$ determinabuntur per aequationes
\[
\begin{aligned}
p&=\frac{rr'(\cos(N-\Pi)-\cos(N'-\Pi))}{r\cos(N-\Pi)-r'\cos(N'-\Pi)}\\[4pt]
e&=\frac{r'-r}{r\cos(N-\Pi)-r'\cos(N'-\Pi)}\mathpunct{\cdot}
\end{aligned}
\]
Denominatorem communem in his formulis reducere licet sub formam $a\cos(A-\Pi)$, ita ut $a$ et $A$ a $\Pi$ sint independentes. Designante scilicet $H$ angulum arbitrarium fit
\[
\begin{aligned}
&r\cos(N-\Pi)-r'\cos(N'-\Pi)\\
&\quad=(r\cos(N-H)-r'\cos(N'-H))\cos(H-\Pi)\\
&\qquad-(r\sin(N-H)-r'\sin(N'-H))\sin(H-\Pi),
\end{aligned}
\]
adeoque $=a\cos(A-\Pi)$, si $a$ et $A$ determinantur per aequationes
\SourceLeaf{110}
\[
\begin{aligned}
r\cos(N-H)-r'\cos(N'-H)&=a\cos(A-H)\\
r\sin(N-H)-r'\sin(N'-H)&=a\sin(A-H).
\end{aligned}
\]
Hoc modo fit
\[
\begin{aligned}
p&=\frac{2rr'\sin\tfrac12(N'-N)\sin(\tfrac12N+\tfrac12N'-\Pi)}{a\cos(A-\Pi)}\\[4pt]
e&=\frac{r'-r}{a\cos(A-\Pi)}\mathpunct{\cdot}
\end{aligned}
\]
Hae formulae imprimis sunt commodae, quoties $p$ et $e$ pro pluribus valoribus ipsius $\Pi$ computandae sunt, manentibus $r$, $r'$, $N$, $N'$. — Quum ad calculum quantitatum auxiliarium $a$, $A$ angulum $H$ ad lubitum assumere liceat, e re erit statuere $H=\tfrac12(N+N')$, quo pacto formulae abeunt in has
\[
\begin{aligned}
(r'-r)\cos\tfrac12(N'-N)&=-a\cos(A-\tfrac12N-\tfrac12N')\\
(r'+r)\sin\tfrac12(N'-N)&=-a\sin(A-\tfrac12N-\tfrac12N').
\end{aligned}
\]
Determinato itaque angulo $A$ per aequationem
\[
\operatorname{tang}(A-\tfrac12N-\tfrac12N')=\frac{r'+r}{r'-r}\operatorname{tang}\tfrac12(N'-N),
\]
statim habetur
\[
e=-\frac{\cos(A-\tfrac12N-\tfrac12N')}{\cos\tfrac12(N'-N)\cos(A-\Pi)}\mathpunct{\cdot}
\]
Calculum logarithmi quantitatis $\dfrac{r'+r}{r'-r}$ per artificium saepius iam explicatum contrahere licebit.
